% Letter 'murica
% 2.23in margins all around gives you the golden ratio between content width
% margins apparently
\documentclass[letterpaper, 12pt]{article}
\usepackage[letterpaper,
	top=1in,
	bottom=1.25in,
	left=1.5in,
	right=1.5in,
]{geometry}

\usepackage[T1]{fontenc}
\usepackage{ETbb}
\let\oldstylenums\textosf

%% Exclusively use indents to demarcate paragraphs
\setlength{\parindent}{0.5in}
%% Don't add any more space than baselineskip after paragraphs
\setlength{\parskip}{0in}

% 135% of the font size is the ideal line spacing
% 1.35 * 12 = 16.20
% 16.20/14.5 (default baselineskip) = 1.116
% \linespread scales \baselineskip ergo the math shenanigans
\linespread{1.117}

\RequirePackage[usenames,dvipsnames]{color}
\definecolor{uvaOrange}{RGB}{229, 114, 0}

\usepackage{fancyhdr}
\pagestyle{fancy}
\renewcommand{\headrulewidth}{0pt}
\renewcommand{\footrulewidth}{0pt}
\fancyhead{}
\fancyfoot{}
\fancyfoot[C]{{\color{uvaOrange}\thepage}}

\usepackage[hidelinks, colorlinks]{hyperref}
% END OF GENERIC DBK_ARTICLE

\usepackage{minted}
\usepackage[page]{appendix}
\usepackage{url}
\usepackage{booktabs}
\usepackage{float}

\usepackage[style=alphabetic]{biblatex}
\addbibresource{ref.bib}

\title{LTL vs SVA generation}
\author{Deebakkarthi Chinnasame Rani}


\usepackage{lipsum}
\begin{document}

%% Title
\maketitle
\thispagestyle{fancy}

%%
\section{Experiment Setup}
I used \cite{sockit} as the DUT. It is a single RTL file but is non-trivial
and implements a real world protocol, namely OneWire \cite{onewire}. The prompts in
Appendix \ref{appendix:prompts} were used. The experiment was repeated 5 times
for both LTL and SVA. I used \texttt{claude-opus-5-5} for all the runs. 

\section{Rudimentary textual metrics}

We see from Table \ref{tab:line_count} that Claude almost always generates more SVA
properties than LTL properties.
\begin{table}[H]
\centering
\begin{tabular}{@{}lll@{}}
\toprule
Run \# & LTL & SVA \\ \midrule
	1      & 46 & \textbf{76}  \\ \midrule
	2      & 43 & \textbf{77}  \\ \midrule
	3      & 71 & \textbf{111}  \\ \midrule
	4      & 63 & \textbf{71}  \\ \midrule
	5      & \textbf{82} & 53  \\ \midrule
Total & 305 & \textbf{388} \\ \bottomrule
\end{tabular}
\caption{Line Count}
\label{tab:line_count}
\end{table}

\subsection{Repetition}
There were 14 \textit{verbatim} repeated properties in LTL files. 10 of them
were in 2 runs, 2 of them were in 4 runs and 2 of them were in all 5 runs.
I had to manually clean up \texttt{"G( "} and \texttt{"G ("}.

There were 21 properties that were repeated verbatim in SVA files. 16 of them
were in 2 runs and 5 of them were in 3 runs. No properties were present in
least 4 runs.

Among the repeated properties only one was the same between LTL and SVA. Another
one was partially related but the other properties had no counterpart in the
other format.

\subsubsection{Logical Equivalence}
I simplified the LTL formulas into \texttt{spot}'s Binary Decision Diagram (BDD). Table
\ref{tab:bdd_equiv} shows that there are lot more repetition now. But there are
still 214 unique BDDs.

\begin{table}[H]
\centering
\begin{tabular}{@{}ll@{}}
\toprule
\#(Repetition) & Count \\ \midrule
1              & 214   \\ \midrule
2              & 30    \\ \midrule
3              & 3     \\ \midrule
4              & 3     \\ \midrule
5              & 2     \\ \bottomrule
\end{tabular}
\caption{Histogram}
\label{tab:bdd_equiv}
\end{table}


\printbibliography
\begin{appendices}
	\section{Prompts}
	\label{appendix:prompts}

	I prompted \texttt{Opus 5.5} to generate the following prompts. I didn't want
	my prompt writing skills to muddle my measurements. I noticed that the prompts
	that it generals makes use of a lot of markdown and is quite verbose and
	precise in its instruction.

	I first asked Claude to read \cite{duret-lutz2022}'s instruction for LTL
	syntax. Specifically I asked it to read \url{https://spot.lre.epita.fr/tl.pdf}
	and summarize them as a list of instructions.

	Similarly Prompt \ref{prompt:sva} was generated by swapping out the word
	\textit{LTL} and replacing it with \textit{SystemVerilog}.

	I had to add explicit instructions to not do anything besides the generation.
	Without this Claude Code tried to run \texttt{yosys} all by
	itself.

	\begin{listing}[h]
		\caption{\texttt{ltl.md}}
		\label{prompt:ltl}
\begin{minted}[breaklines]{markdown}
Read `sockit_owm.v` and write linear temporal properties of its behavior to `sockit_owm.ltl`, one formula per line. Reading that file and writing this one are the only actions permitted; do not take or attempt any other action.

Use this syntax:

- Constants: `1`, `0`.
- Atomic propositions: lowercase identifiers. Double-quote any name that starts with uppercase F, G, or X, matches a keyword (F G M R U V W X xor true false), or is an expression, e.g. `"cnt==3"`, `"Fsm"`.
- Boolean operators: `!` `&` `|` `->` `<->` `xor`.
- Temporal operators: `X` (next), `F` (eventually), `G` (always), `U` (until), `W` (weak until), `R` (release), `M` (strong release).
- Bounded forms: `X[n] f`, `F[n:m] f`, `G[n:m] f`.
- Precedence, tightest first: `!`, then `X`, then `F`/`G`, then `U W R M` (right-associative), then `&`, `|`, `xor`, and finally `->` and `<->` (right-associative). Add parentheses whenever grouping is unclear.
\end{minted}
	\end{listing}

	\begin{listing}[h]
		\caption{\texttt{sva.md}}
		\label{prompt:sva}
\begin{minted}[breaklines]{markdown}
Read `sockit_owm.v` and write SystemVerilog Assertions of its behavior to `sockit_owm.sva`, one assertion per line. Reading that file and writing this one are the only actions permitted; do not take or attempt any other action.
\end{minted}
	\end{listing}

\end{appendices}
\end{document}
